\section{What the literature says about each determinant}\label{sec:lit}

\textit{Population and immigration.} In spatial equilibrium, local rents capitalise the demand of residents who choose where to live \citep{roback1982}. Immigrant inflows provide a source of local demand shocks that can be predicted from earlier settlement patterns. With that strategy, \citet{saiz2007} estimates that an inflow equal to 1\,\% of a US metropolitan area's population raises rents by about 1\,\%, and \citet{sa2015} finds negative effects on house prices in England, explained by the out-migration of higher-income natives. For Spain, \citet{gonzalez2013} relate the immigration boom of 2000--2010 to house prices and construction, and \citet{sanchis2023} decomposes the price response by location. Shift-share instruments of this kind rely either on exogenous shares or on many exogenous shocks \citep{goldsmith2020,borusyak2022,adao2019}, and they can confound short- and long-run responses when settlement patterns are persistent.

\textit{Income, credit and the user cost.} Rents also capitalise local incomes and amenities \citep{roback1982,gyourko2013}, and the choice between renting and owning depends on the user cost of owner-occupation, which moves with interest rates, credit standards and expected price growth \citep{poterba1984,himmelberg2005,sinai2005}. These forces are largely common to all municipalities in a country, which is why the cross-sectional literature on local rents seldom identifies them. For Spain, \citet{khametshin2024} document that the rise in rental demand since 2015 is concentrated among young and foreign-born households and in large urban and tourist areas, and attribute part of it to prudent mortgage lending; they explicitly describe the quantification of each factor as an open question.

\textit{Supply.} Where supply is inelastic, demand shocks are capitalised into prices rather than accommodated by construction. Inelasticity reflects geography \citep{saiz2010} and regulation \citep{glaeser2018,hilber2016,molloy2020}, and housing durability makes supply asymmetric \citep{glaeser2005}. Constraints in productive cities have aggregate costs \citep{hsieh2019}. \citet{khametshin2024} show that rents grew more in large Spanish municipalities with less land available for building.

\textit{Tourist dwellings.} Home-sharing platforms raise long-term rents by moving dwellings from residents to visitors. \citet{barron2021} formalise this reallocation and estimate it with a shift-share instrument for the US; \citet{garcialopez2020} find that in Barcelona neighbourhoods with high Airbnb activity rents rose by several percentage points; \citet{koster2021} use regulatory borders in Los Angeles; and \citet{franco2021} and \citet{duso2024} report evidence for Portugal and Berlin. \citet{batalha2022} exploit the collapse of tourism in 2020 in Lisbon, which pushed short-term-rental dwellings back into the residential market and lowered rents. \citet{chaves2026} finds that Airbnb raised house prices and reduced residential inflows in Madrid neighbourhoods, and \citet{almagro2025} show how tourism reshapes neighbourhood amenities and sorting. Regulation of short-term rentals has its own effects \citep{valentin2021}.

\textit{Rent regulation.} Rent control lowers rents in covered units but reduces the supply of rental housing and can raise rents in uncovered segments \citep{autor2014,diamond2019}; \citet{kholodilin2024} reviews the evidence. In Catalonia, the 2020--2022 rent-containment law reduced rents and the number of leases signed \citep{jofre2023,monras2023}, and a first evaluation of the 2024 tensioned-zone cap finds a reduction in the number of new leases and a less robust effect on rent growth \citep{perez2026}.

\textit{Institutional investors.} In the United States, large buy-to-rent investors bought foreclosed single-family homes after 2008 \citep{mills2019}; the evidence on their effect on rents is mixed, with higher rents and fees in some markets \citep{gurun2023} and price stabilisation in others \citep{garriga2023}. In Spain, sociological and geographical work has documented the acquisition of rental portfolios by global funds and their role in tourist rentals \citep{janoschka2020,cocola2025}, but there is no causal estimate of their effect on rents. According to tax data compiled by \citet{khametshin2024}, companies owned 8.1\,\% of the dwellings let as main residences in 2021, ranging from 2.6\,\% to 11.1\,\% across regions.

\textit{Contribution of this paper.} Existing work studies one determinant at a time, usually in one city. This paper estimates several channels within a common framework, for all Spanish municipalities with an official rent index, distinguishes the stock of leases from new contracts, and states for each channel what the design identifies.
